\documentclass{guia}
\grado{3}
\nivel{Secundaria}
\cicloescolar{2026-2027}
\materia{Matemáticas}
\guia{1}
\unidad{1}
\title{Cálculos Numéricos}
\aprendizajes{
      \item Convierte fracciones decimales a notación decimal y viceversa. Aproxima algunas fracciones no decimales usando la notación decimal. 
      }
\author{Melchor Pinto, JC}
\begin{document}
\INFO%
% \begin{multicols}{2}
% 	\tableofcontents
% \end{multicols}
% \newpage
\begin{questions}
      \section{Suma de Números}
      \questionboxed{Realiza las siguientes sumas:

	\begin{multicols}{4}
		\begin{parts}
			\part $9 + 8=$ \fillin[17][0cm]\\[0.2cm]
			\part $5 + 4=$ \fillin[9][0cm]\\[0.2cm]

			\part
			\ifprintanswers{\opadd[hfactor=decimal,resultstyle=\color{red},carryadd=true]{17}{18}}
			\else{\opadd[hfactor=decimal,resultstyle=\color{white},carryadd=false]{17}{18}}\fi\\[0.2cm]

\columnbreak%

			\part $1 + 1=$ \fillin[2][0cm]\\[0.2cm]
			\part $5 + 7=$ \fillin[12][0cm]\\[0.2cm]

			\part
			\ifprintanswers{\opadd[hfactor=decimal,resultstyle=\color{red},carryadd=true]{15}{9}}
			\else{\opadd[hfactor=decimal,resultstyle=\color{white},carryadd=false]{15}{9}}\fi\\[0.2cm]
			
			\columnbreak%
			% \part
			% \ifprintanswers{\opadd[hfactor=decimal,resultstyle=\color{red},carryadd=true]{26}{9}}
			% \else{\opadd[hfactor=decimal,resultstyle=\color{white},carryadd=false]{26}{9}\\[0.5cm]}\fi\\[0.2cm]

			% \part
			% \ifprintanswers{\opadd[hfactor=decimal,resultstyle=\color{red},carryadd=true]{27}{12}}
			% \else{\opadd[hfactor=decimal,resultstyle=\color{white},carryadd=false]{27}{12}}\fi\\[0.2cm]

			\part $0 + 7=$ \fillin[7][0cm]\\[0.2cm]
			\part $8 + 7=$ \fillin[15][0cm]\\[0.2cm]

			\part
			\ifprintanswers{\opadd[hfactor=decimal,resultstyle=\color{red},carryadd=true]{10}{9}}
			\else{\opadd[hfactor=decimal,resultstyle=\color{white},carryadd=false]{10}{9}}\fi\\[0.2cm]
		
			\columnbreak%
			% \part
			% \ifprintanswers{\opadd[hfactor=decimal,resultstyle=\color{red},carryadd=true]{37}{8}}
			% \else{\opadd[hfactor=decimal,resultstyle=\color{white},carryadd=false]{37}{8}\\[0.5cm]}\fi\\[0.2cm]

			% \part
			% \ifprintanswers{\opadd[hfactor=decimal,resultstyle=\color{red},carryadd=true]{48}{14}}
			% \else{\opadd[hfactor=decimal,resultstyle=\color{white},carryadd=false]{48}{14}}\fi\\[0.2cm]

			\part $1 + 9=$ \fillin[10][0cm]\\[0.2cm]
			\part $4 + 9=$ \fillin[13][0cm]\\[0.2cm]

			\part
			\ifprintanswers{\opadd[hfactor=decimal,resultstyle=\color{red},carryadd=true]{21}{19}}
			\else{\opadd[hfactor=decimal,resultstyle=\color{white},carryadd=false]{21}{19}}\fi\\[0.2cm]
			% \part
			% \ifprintanswers{\opadd[hfactor=decimal,resultstyle=\color{red},carryadd=true]{44}{25}}
			% \else{\opadd[hfactor=decimal,resultstyle=\color{white},carryadd=false]{44}{25}\\[0.5cm]}\fi\\[0.2cm]

			% \part
			% \ifprintanswers{\opadd[hfactor=decimal,resultstyle=\color{red},carryadd=true]{82}{38}}
			% \else{\opadd[hfactor=decimal,resultstyle=\color{white},carryadd=false]{82}{38}}\fi\\[0.2cm]
		\end{parts}
	\end{multicols}
}

	\questionboxed{Realiza las siguientes sumas:

		\begin{multicols}{4}
			\begin{parts}
				\part \ifprintanswers{\large  \quad   \opadd[hfactor=decimal,resultstyle=\color{red},carryadd=true,carrysub=false]{37854}{18581} }
				\else{          \large  \quad  \opadd[hfactor=decimal,resultstyle=\color{white},carryadd=false,carrysub=false]{37854}{18581} }
				\fi
				\part \ifprintanswers{\large  \quad   \opadd[hfactor=decimal,resultstyle=\color{red},carryadd=true,carrysub=false]{3234}{24156} }
				\else{          \large  \quad  \opadd[hfactor=decimal,resultstyle=\color{white},carryadd=false,carrysub=false]{3234}{24156} }
				\fi
				\part \ifprintanswers{\large  \quad   \opadd[hfactor=decimal,resultstyle=\color{red},carryadd=true,carrysub=false]{30985}{19562} }
				\else{          \large  \quad  \opadd[hfactor=decimal,resultstyle=\color{white},carryadd=false,carrysub=false]{30985}{19562} }
				\fi
				\part \ifprintanswers{\large  \quad   \opadd[hfactor=decimal,resultstyle=\color{red},carryadd=true,carrysub=false]{2849}{2415} }
				\else{          \large  \quad  \opadd[hfactor=decimal,resultstyle=\color{white},carryadd=false,carrysub=false]{2849}{2415} }
				\fi
				\part \ifprintanswers{\large  \quad   \opadd[hfactor=decimal,resultstyle=\color{red},carryadd=true,carrysub=false]{31085}{19001} }
				\else{          \large  \quad  \opadd[hfactor=decimal,resultstyle=\color{white},carryadd=false,carrysub=false]{31085}{19001} }
				\fi
				\part \ifprintanswers{\large  \quad   \opadd[hfactor=decimal,resultstyle=\color{red},carryadd=true,carrysub=false]{35701}{25484} }
				\else{          \large  \quad  \opadd[hfactor=decimal,resultstyle=\color{white},carryadd=false,carrysub=false]{35701}{25484} }
				\fi
				\part \ifprintanswers{\large  \quad   \opadd[hfactor=decimal,resultstyle=\color{red},carryadd=true,carrysub=false]{45668}{19624} }
				\else{          \large  \quad  \opadd[hfactor=decimal,resultstyle=\color{white},carryadd=false,carrysub=false]{45668}{19624} }
				\fi
				\part \ifprintanswers{\large  \quad   \opadd[hfactor=decimal,resultstyle=\color{red},carryadd=true,carrysub=false]{58718}{3652} }
				\else{          \large  \quad  \opadd[hfactor=decimal,resultstyle=\color{white},carryadd=false,carrysub=false]{58718}{3652} }
				\fi
			\end{parts}
		\end{multicols}
	}

      

      \section{Resta de Números}
      \questionboxed[3]{Realiza las siguientes restas:

	\begin{multicols}{4}
		\begin{parts}
			\part $9 - 3=$ \fillin[6][0cm]\\[0.2cm]
			% \part $15 - \fillin[8][0.5cm]= 7$\\[0.2cm]

			\part
			\ifprintanswers{\opsub[hfactor=decimal,resultstyle=\color{red},carrysub=true]{17}{6}}
			\else{\opsub[hfactor=decimal,resultstyle=\color{white},carrysub=false]{17}{6}}\fi\\[0.2cm]

			\part
			\ifprintanswers{\opsub[hfactor=decimal,resultstyle=\color{red},carrysub=true]{15}{4}}
			\else{\opsub[hfactor=decimal,resultstyle=\color{white},carrysub=false]{15}{4}}\fi\\[0.2cm]

			% \part
			% \ifprintanswers{\opsub[hfactor=decimal,resultstyle=\color{red},carrysub=true]{34}{22}}
			% \else{\opsub[hfactor=decimal,resultstyle=\color{white},carrysub=false]{34}{22}}\fi\\[0.2cm]

			\part $7 - 4=$ \fillin[3][0cm]\\[0.2cm]
			% \part $12 - \fillin[7][0.5cm]= 5$\\[0.2cm]

			\part
			\ifprintanswers{\opsub[hfactor=decimal,resultstyle=\color{red},carrysub=true]{27}{5}}
			\else{\opsub[hfactor=decimal,resultstyle=\color{white},carrysub=false]{27}{5}}\fi\\[0.2cm]

			\part
			\ifprintanswers{\opsub[hfactor=decimal,resultstyle=\color{red},carrysub=true]{14}{11}}
			\else{\opsub[hfactor=decimal,resultstyle=\color{white},carrysub=false]{14}{11}}\fi\\[0.2cm]

			% \part
			% \ifprintanswers{\opsub[hfactor=decimal,resultstyle=\color{red},carrysub=true]{48}{29}}
			% \else{\opsub[hfactor=decimal,resultstyle=\color{white},carrysub=false]{48}{29}}\fi\\[0.2cm]

			\part $8 - 8=$ \fillin[0][0cm]\\[0.2cm]
			% \part $18 - \fillin[14][0.5cm]= 4$\\[0.2cm]

			\part
			\ifprintanswers{\opsub[hfactor=decimal,resultstyle=\color{red},carrysub=true]{8}{5}}
			\else{\opsub[hfactor=decimal,resultstyle=\color{white},carrysub=false]{8}{5}}\fi\\[0.2cm]

			\part
			\ifprintanswers{\opsub[hfactor=decimal,resultstyle=\color{red},carrysub=true]{16}{9}}
			\else{\opsub[hfactor=decimal,resultstyle=\color{white},carrysub=false]{16}{9}}\fi\\[0.2cm]

			% \part
			% \ifprintanswers{\opsub[hfactor=decimal,resultstyle=\color{red},carrysub=true]{57}{43}}
			% \else{\opsub[hfactor=decimal,resultstyle=\color{white},carrysub=false]{57}{43}}\fi\\[0.2cm]

			\part $11 - 4=$ \fillin[7][0cm]\\[0.2cm]
			% \part $25 - \fillin[20][0.5cm]= 5$\\[0.2cm]

			\part
			\ifprintanswers{\opsub[hfactor=decimal,resultstyle=\color{red},carrysub=true]{17}{5}}
			\else{\opsub[hfactor=decimal,resultstyle=\color{white},carrysub=false]{17}{5}}\fi\\[0.2cm]

			\part
			\ifprintanswers{\opsub[hfactor=decimal,resultstyle=\color{red},carrysub=true]{10}{8}}
			\else{\opsub[hfactor=decimal,resultstyle=\color{white},carrysub=false]{10}{8}}\fi\\[0.2cm]

			% \part
			% \ifprintanswers{\opsub[hfactor=decimal,resultstyle=\color{red},carrysub=true]{63}{18}}
			% \else{\opsub[hfactor=decimal,resultstyle=\color{white},carrysub=false]{63}{18}}\fi\\[0.2cm]
		\end{parts}
	\end{multicols}
}


	\questionboxed[12]{Realiza las siguientes restas:

		\begin{multicols}{4}
			\begin{parts}
				\part \ifprintanswers{\large  \quad   \opsub[hfactor=decimal,resultstyle=\color{red},carryadd=true,carrysub=true]{4000}{2267} }
				\else{          \large  \quad  \opsub[hfactor=decimal,resultstyle=\color{white},carryadd=false,carrysub=false]{4000}{2267} }
				\fi
				\part \ifprintanswers{\large  \quad   \opsub[hfactor=decimal,resultstyle=\color{red},carryadd=true,carrysub=true]{800}{744} }
				\else{          \large  \quad  \opsub[hfactor=decimal,resultstyle=\color{white},carryadd=false,carrysub=false]{800}{744} }
				\fi
				\part \ifprintanswers{\large  \quad   \opsub[hfactor=decimal,resultstyle=\color{red},carryadd=true,carrysub=true]{3500}{308} }
				\else{          \large  \quad  \opsub[hfactor=decimal,resultstyle=\color{white},carryadd=false,carrysub=false]{3500}{308} }
				\fi
				\part \ifprintanswers{\large  \quad   \opsub[hfactor=decimal,resultstyle=\color{red},carryadd=true,carrysub=true]{3000}{189} }
				\else{          \large  \quad  \opsub[hfactor=decimal,resultstyle=\color{white},carryadd=false,carrysub=false]{3000}{189} }
				\fi
				\part \ifprintanswers{\large  \quad   \opsub[hfactor=decimal,resultstyle=\color{red},carryadd=true,carrysub=true]{1200}{966} }
				\else{          \large  \quad  \opsub[hfactor=decimal,resultstyle=\color{white},carryadd=false,carrysub=false]{1200}{966} }
				\fi
				\part \ifprintanswers{\large  \quad   \opsub[hfactor=decimal,resultstyle=\color{red},carryadd=true,carrysub=true]{3300}{2117} }
				\else{          \large  \quad  \opsub[hfactor=decimal,resultstyle=\color{white},carryadd=false,carrysub=false]{3300}{2117} }
				\fi
				\part \ifprintanswers{\large  \quad   \opsub[hfactor=decimal,resultstyle=\color{red},carryadd=true,carrysub=true]{2000}{1251} }
				\else{          \large  \quad  \opsub[hfactor=decimal,resultstyle=\color{white},carryadd=false,carrysub=false]{2000}{1251} }
				\fi
				\part \ifprintanswers{\large  \quad   \opsub[hfactor=decimal,resultstyle=\color{red},carryadd=true,carrysub=true]{2400}{2023} }
				\else{          \large  \quad  \opsub[hfactor=decimal,resultstyle=\color{white},carryadd=false,carrysub=false]{2400}{2023} }
				\fi
				\part \ifprintanswers{\large  \quad   \opsub[hfactor=decimal,resultstyle=\color{red},carryadd=true,carrysub=false]{2400}{211} }
				\else{          \large  \quad  \opsub[hfactor=decimal,resultstyle=\color{white},carryadd=false,carrysub=false]{2400}{211} }
				\fi
				\part \ifprintanswers{\large  \quad   \opsub[hfactor=decimal,resultstyle=\color{red},carryadd=true,carrysub=false]{1500}{1044} }
				\else{          \large  \quad  \opsub[hfactor=decimal,resultstyle=\color{white},carryadd=false,carrysub=false]{1500}{1044} }
				\fi
				\part \ifprintanswers{\large  \quad   \opsub[hfactor=decimal,resultstyle=\color{red},carryadd=true,carrysub=false]{2000}{1105} }
				\else{          \large  \quad  \opsub[hfactor=decimal,resultstyle=\color{white},carryadd=false,carrysub=false]{2000}{1105} }
				\fi
				\part \ifprintanswers{\large  \quad   \opsub[hfactor=decimal,resultstyle=\color{red},carryadd=true,carrysub=false]{1600}{669} }
				\else{          \large  \quad  \opsub[hfactor=decimal,resultstyle=\color{white},carryadd=false,carrysub=false]{1600}{669} }
				\fi
			\end{parts}
		\end{multicols}
	}

      \section{Multiplicación de Números}
     \questionboxed[15]{Reponde las siguientes tablas de multiplicar:

		\begin{multicols}{4}
			\begin{parts}\Large
				\part $7 \times 6=$ \fillin[$42$][0cm]
				\part $9 \times 7=$ \fillin[$63$][0cm]
				\part $6 \times 8=$ \fillin[$48$][0cm]
				\part $6 \times 9=$ \fillin[$54$][0cm]
				\part $4 \times 4=$ \fillin[$16$][0cm]
				\part $7 \times 7=$ \fillin[$49$][0cm]
				\part $7 \times 5=$ \fillin[$35$][0cm]
				\part $5 \times 4=$ \fillin[$20$][0cm]
				\part $8 \times 7=$ \fillin[$56$][0cm]
				\part $3 \times 6=$ \fillin[$18$][0cm]
				\part $5 \times 9=$ \fillin[$45$][0cm]
				\part $5 \times 6=$ \fillin[$30$][0cm]
				\part $3 \times 8=$ \fillin[$24$][0cm]
				\part $2 \times 9=$ \fillin[$18$][0cm]
				\part $2 \times 7=$ \fillin[$14$][0cm]
				\part $4 \times 7=$ \fillin[$28$][0cm]
				\part $3 \times 9=$ \fillin[$27$][0cm]
				\part $5 \times 8=$ \fillin[$40$][0cm]
				\part $6 \times 6=$ \fillin[$36$][0cm]
				\part $7 \times 8=$ \fillin[$56$][0cm]
				\part $9 \times 4=$ \fillin[$36$][0cm]
				\part $8 \times 8=$ \fillin[$64$][0cm]
				\part $9 \times 9=$ \fillin[$81$][0cm]
				\part $6 \times 7=$ \fillin[$42$][0cm]
			\end{parts}
		\end{multicols}
	}

	\questionboxed[15]{Completa las siguientes tablas de multiplicar:

		\begin{multicols}{4}
			\begin{parts}\Large
				\part $\fillin[9][0.5cm] \times 9= 81$
				\part $4 \times \fillin[9][0.5cm]= 36$
				\part $\fillin[7][0.5cm] \times 4= 28$
				\part $\fillin[9][0.5cm] \times 3= 27$
				\part $8 \times \fillin[5][0.5cm]= 40$
				\part $\fillin[6][0.5cm] \times 4= 24$
				\part $7 \times \fillin[7][0.5cm]= 49$
				\part $\fillin[8][0.5cm] \times 3= 24$
				\part $9 \times \fillin[8][0.5cm]= 72$
				\part $10 \times \fillin[10][0.5cm]=100$
				\part $5 \times \fillin[7][0.5cm]=35$
				\part $6 \times \fillin[7][0.5cm]= 42$	
				\part $\fillin[1][0.5cm] \times 9= 9$
				\part $\fillin[6][0.5cm] \times 5= 30$
				\part $\fillin[7][0.5cm] \times 3= 21$
				
				\part $9 \times \fillin[10][0.5cm]=90$

				\part $\fillin[7][0.5cm] \times 8= 56$
				\part $5 \times \fillin[10][0.5cm]=50$
				\part $4 \times \fillin[6][0.5cm]=24$
				\part $1 \times \fillin[7][0.5cm]= 7$	
				\part $\fillin[9][0.5cm] \times 5= 45$
				\part $\fillin[6][0.5cm] \times 6= 36$
				\part $\fillin[8][0.5cm] \times 8= 64$
				\part $\fillin[6][0.5cm] \times 9= 54$
				\part	$3 \times \fillin[10][0.5cm]=30$
				\part $4 \times \fillin[8][0.5cm]=32$
				\part $6 \times \fillin[0][0.5cm]= 0$				
			\end{parts}
		\end{multicols}
	}

      \questionboxed[6]{Realiza las siguientes multiplicaciones:

		\begin{multicols}{3}
			\begin{parts}
				\part \ifprintanswers{\large  \quad   \opmul[hfactor=decimal,resultstyle=\color{red},displayintermediary=None]{314}{2} }
				\else{          \large  \quad  \opmul[hfactor=decimal,resultstyle=\color{white},displayintermediary=None]{314}{2} }
				\fi
				\part \ifprintanswers{\large  \quad   \opmul[hfactor=decimal,resultstyle=\color{red},displayintermediary=None]{283}{4} }
				\else{          \large  \quad  \opmul[hfactor=decimal,resultstyle=\color{white},displayintermediary=None]{283}{4} }
				\fi
				\part \ifprintanswers{\large  \quad   \opmul[hfactor=decimal,resultstyle=\color{red},displayintermediary=None]{2781}{5} }
				\else{          \large  \quad  \opmul[hfactor=decimal,resultstyle=\color{white},displayintermediary=None]{2781}{5} }
				\fi
				\part \ifprintanswers{\large  \quad   \opmul[hfactor=decimal,resultstyle=\color{red},displayintermediary=None]{4914}{6} }
				\else{          \large  \quad  \opmul[hfactor=decimal,resultstyle=\color{white},displayintermediary=None]{4914}{6} }
				\fi
				\part \ifprintanswers{\large  \quad   \opmul[hfactor=decimal,resultstyle=\color{red},displayintermediary=None]{255}{24} }
				\else{          \large  \quad  \opmul[hfactor=decimal,resultstyle=\color{white},displayintermediary=None]{255}{24} }
				\fi
				\part \ifprintanswers{\large  \quad   \opmul[hfactor=decimal,resultstyle=\color{red},displayintermediary=None]{3533}{29} }
				\else{          \large  \quad  \opmul[hfactor=decimal,resultstyle=\color{white},displayintermediary=None]{3533}{29} }
				\fi
			\end{parts}
		\end{multicols}
	}
      
      \section{División de Números}
\questionboxed[8]{Realiza las siguientes divisiones:

		\begin{multicols}{4}
			\begin{parts}
				\part \ifprintanswers{\large\opidiv{123}{6}} \\[0.5em]
				\else{          \Large  \quad $6 \overline{) \ 123\ }$} \\[2em]
				\fi

				\part \ifprintanswers{\large\opidiv{200}{3}} \\[0.5em]
				\else{          \Large  \quad $3 \overline{) \ 200\ }$} \\[2em]
				\fi

				\part \ifprintanswers{\large\opidiv{399}{8}} \\[0.5em]
				\else{          \Large  \quad $8 \overline{) \ 399\ }$} \\[2em]
				\fi

				\part \ifprintanswers{\large\opidiv{193}{7}} \\[0.5em]
				\else{          \Large  \quad $7 \overline{) \ 193\ }$} \\[2em]
				\fi

				\part \ifprintanswers{\large\opidiv{283}{6}} \\[0.5em]
				\else{          \Large  \quad $6 \overline{) \ 283\ }$} \\[2em]
				\fi

				\part \ifprintanswers{\large\opidiv{432}{9}} \\[0.5em]
				\else{          \Large  \quad $9 \overline{) \ 432\ }$} \\[2em]
				\fi

				\part \ifprintanswers{\large\opidiv{644}{8}} \\[0.5em]
				\else{          \Large  \quad $8 \overline{) \ 644\ }$} \\[2em]
				\fi

				\part \ifprintanswers{\large\opidiv{656}{7}} \\[0.5em]
				\else{          \Large  \quad $7 \overline{) \ 656\ }$} \\[2em]
				\fi
			\end{parts}
		\end{multicols}
	}
    
      \section{Resolución de Problemas}
\questionboxed[2]{Resuelve los siguientes problemas sobre sumas y restas:

		\begin{multicols}{2}
			\begin{parts}
				\part El total de mis compras es de 315 pesos, ¿cuánto dinero recibiré de cambio si pago con un billete de 500 pesos?

				\begin{solutionbox}{1cm}
					\opsub[style=text]{500}{315}
				\end{solutionbox}

				\part Luis tiene ahorrado 257 pesos, si su abuelo le regala 360 pesos más, ¿cuánto dinero tiene en total Luis?

				\begin{solutionbox}{1cm}
					\opadd[style=text]{257}{360}
				\end{solutionbox}

				\part Jorge está armando un rompecabezas de 500 piezas, si ha puesto 233 piezas, ¿cuántas piezas le faltan por poner a Jorge?

				\begin{solutionbox}{1cm}
					\opsub[style=text]{500}{233}
				\end{solutionbox}

				\part Carlos mide 183 centímetros y es 8 centímetros más alto que Julio, ¿cuántos centímetros mide Julio?

				\begin{solutionbox}{1cm}
					\opsub[style=text]{183}{8}
				\end{solutionbox}
			\end{parts}
		\end{multicols}
	}

      \questionboxed[5]{Resuelve los siguientes problemas sobre multiplicaciones:

		\begin{multicols}{2}
			\begin{parts}
				\part Una escuela tiene 6 salones, si cada salón tiene 25 alumnos. ¿Cuántos alumnos tiene en total la escuela?

				\begin{solutionbox}{1cm}
					\opmul[style=text]{6}{25}
				\end{solutionbox}

				\part Una cubeta de pintura cuesta 2345 pesos, ¿cuánto se pagará por 3 cubetas de pintura?

				\begin{solutionbox}{1cm}
					\opmul[style=text]{3}{2345}
				\end{solutionbox}

				\part Una secretaria puede escribir 36 palabras por minuto si continua con este ritmo, ¿cuántas palabras puede escribir en 12 minutos?

				\begin{solutionbox}{1cm}
					\opmul[style=text]{36}{12}
				\end{solutionbox}

				\part Cristina compró 5 cajas de leche de soya, si cada caja tiene 12 envases de leche, ¿cuántos envases de leche compró Cristina?

				\begin{solutionbox}{1cm}
					\opmul[style=text]{5}{12}
				\end{solutionbox}

				\part Mariana fue a la frutería y compró 3 kilogramos de uvas, si el kilogramo cuesta 84 pesos. ¿Cuánto pagó en total Mariana?

				\begin{solutionbox}{1cm}
					\opmul[style=text]{3}{84}
				\end{solutionbox}

				\part Laura compró 28 paquetes de galletas, si cada paquete tiene 18 galletas. ¿Cuántas galletas tiene en total Laura?

				\begin{solutionbox}{1cm}
					\opmul[style=text]{28}{18}
				\end{solutionbox}

			\end{parts}
		\end{multicols}
	}


\end{questions}
\end{document}